\section{Discussion and Open Questions}
\label{sec:discussion}

\subsection{What this paper establishes}

\paragraph{Direct $\volR$ measurement is privacy-incompatible.}
The framework's commitments
in~\citep{lasser2026exo,lasser2026phase,lasser2026wamura} rule out
the panopticon-style observation that direct measurement would
require.  This is a structural finding, not a practical limitation.

\paragraph{Revealed sacrifice is a privacy-minimal surrogate.}
The sacrifice channel satisfies five privacy properties
(Theorem~\ref{thm:privacy}) and composes with the commitment layer
of~\citep{lasser2026exo}
(Propositions~\ref{prop:commit_preserves}
and~\ref{prop:zk_aggregation}) to produce a privacy-minimal
observation mechanism.

\paragraph{The B-to-C gap is lower-bounded by characterized trade
data.}
The aggregate lower bound
(Theorem~\ref{thm:aggregate_bound}) converts the B-to-C gap from
``unknown divergence'' into ``lower-bounded divergence.''  The
bound tightens monotonically with trade coverage
(Proposition~\ref{prop:monotone}).

\paragraph{Two sacrifice channels cover complementary territory.}
The money channel captures wealth-correlated capabilities; the
time channel captures identity-constitutive capabilities.
Together they cover more of $\volR$ than either alone
(Corollary~\ref{cor:joint_coverage}).

\paragraph{$\volR$ is not self-balancing (retained finding).}
$\volR$ itself fails M6 structurally
(Theorem~\ref{thm:axiom_inheritance}, failure~(a)), so $\volR$ is
a diagnostic, not an objective.  This structural finding is
independent of the observation mechanism.

\paragraph{Wireheading is detectable from trade data.}
Trade-flow concentration
(Proposition~\ref{prop:wireheading_hhi}) provides a tractable
wireheading diagnostic computable from public committed sacrifice
events
(Proposition~\ref{prop:third_party_wireheading}).

\paragraph{Need-sufficiency resolves the polarity problem.}
The sufficiency threshold
(Section~\ref{sec:need_sufficiency}) identifies the structural
boundary where sacrifice polarity flips.  Below-sufficiency
sacrifices are involuntary and correctly filtered by S1.  The
need-access cost diagnostic
(Section~\ref{sec:gap_decomposition},
Definitions~\ref{def:nac}--\ref{def:nse}) tracks
need-satisfaction efficiency as a companion metric.

\paragraph{Complementarity with controlled relaxation.}
Together, the present paper
and~\citep{lasser2026wamura} cover the risk and value dimensions
of the restriction criterion.  Neither is sufficient alone.

\subsection{Connection to the capability approach}

The distinction between basic capabilities (functionings required
for a minimally decent life) and complex capabilities (freedoms to
choose among valuable life-options) in~\citet{sen1999development}
and~\citet{nussbaum2011creating} maps to the
below-sufficiency/above-sufficiency partition
(Section~\ref{sec:need_sufficiency}).  What the framework adds
beyond that tradition:

\begin{enumerate}
\item \textbf{Downstream-safety condition}
  (Definition~\ref{def:downstream_safe}) that makes ``sufficient''
  structurally rigorous rather than normatively asserted.
\item \textbf{Polarity-correct benchmark discipline}
  (Definition~\ref{def:polarity_benchmark}) that makes
  $\volL$-maximization align with the agent's interest on both
  sides of the threshold.
\item \textbf{Revealed-sacrifice transition} (threshold-setting
  mechanism~(c), Section~\ref{sec:threshold_setting}) as an
  empirical method for identifying the sufficiency threshold
  without normative declaration.
\end{enumerate}

\subsection{Limitations}

\paragraph{Assumption load.}
The full per-capability refinement (Theorem~\ref{thm:aggregate_bound}
Part~B) requires six assumptions (S0--S5).  S0 (calibration) and
S1 (free choice) are the most empirically demanding, and neither
is ZK-verifiable (Remark~\ref{rem:zk_limits}).  Part~A (S0--S2
only) is the more defensible result; Part~B should be understood
as the refinement available under stronger conditions.

\paragraph{B-to-C gap inheritance.}
The damage bound of~\citep[Theorem~1]{lasser2026wamura} bounds
contraction of $\volL$, not of the experiential target directly.
The present paper provides a lower bound on $\volR$ as a
diagnostic, but the diagnostic is itself subject to the B-to-C gap
on the benchmarked side: the sacrifice signal $\Delta\volL(X)$ is
only as accurate as $\volL$'s measurement of the benchmarked
capability.

\paragraph{Residual class irreducibility.}
Capabilities in $\ResS$ (Definition~\ref{def:residual_class}) have
zero sacrifice evidence by construction and cannot gain it
regardless of observation time.  The residual class is the
irreducible component of the gap decomposition.

\subsection{Open questions}

\begin{enumerate}

\item \label{oq:hedonic}
\textbf{Hedonic disaggregation of bundles.}
Theorem~\ref{thm:aggregate_bound} Part~B requires event-local
decomposition coefficients $\alpha_{n,c}$.  Classical hedonic
regression~\citep{rosen1974hedonic} handles the money channel.
The time-sacrifice channel needs an analogous framework: how does
one decompose a block of time spent on a complex activity
(parenting, which involves teaching, emotional regulation,
physical care, relationship-building) into per-capability
contributions?

\item \label{oq:nonstationarity}
\textbf{Non-stationarity of revealed preferences.}
Preferences shift over time.  The lower bound from past trades may
overstate $\volR$ at present.  An EWMA decay paralleling the
risk-trust dynamics of~\citep[Definition~4]{lasser2026scm} seems
natural but requires empirical calibration.

\item \label{oq:coercion}
\textbf{Coerced sacrifice.}
Assumption S1 (free choice) is load-bearing.  Under duress, the
agent's trade does not reveal preferences.  Filtering coerced
sacrifice events may require commitment-based attestation of
choice sets~\citep{lasser2026exo}.

\item \label{oq:market_failure}
\textbf{Market-failure cases.}
In contexts without functioning markets (pre-capitalist societies,
closed economies, intra-family sharing), the money-sacrifice
channel is quiet even though $\volR$ activity is real.  Does the
time-sacrifice channel alone suffice, or does the framework need
a non-market sacrifice variant?

\item \label{oq:commitment_infra}
\textbf{Commitment infrastructure at scale.}
A fully committed trade ledger at the scale of a modern economy is
an infrastructure problem.  The privacy guarantee
(Section~\ref{sec:commitment}) is conditional on the commitment
infrastructure existing, which is the remit
of~\citep{lasser2026exo}.

\item \label{oq:calibration}
\textbf{Population-level calibration (S0).}
Assumption S0 requires that individual agents do not overvalue
unbenchmarked bundles.  A robust version of
Theorem~\ref{thm:aggregate_bound} would replace S0 with a weaker
population-level condition:
$E[\Ui(Y)] \leq \volR(Y) + \epsilon$ for some bounded bias
$\epsilon$.

\item \label{oq:threshold_sensitivity}
\textbf{Threshold sensitivity.}
How sensitive are the alignment properties to the exact placement
of sufficiency thresholds?  If the threshold is set too low,
need-costs in the gap between threshold and true sufficiency are
misclassified as want-pursuit events.  If set too high, the
framework over-constrains the need.

\item \label{oq:dynamic_thresholds}
\textbf{Dynamic thresholds.}
Sufficiency levels may shift over time as technology and social
context evolve.  The framework needs a mechanism for updating
thresholds, connecting to the idea of controlled relaxation
applied to sufficiency standards themselves.

\item \label{oq:cross_cultural}
\textbf{Cross-cultural variation.}
Different cultures may set sufficiency thresholds differently.
Can the framework maintain culturally-indexed thresholds while
preserving a coherent aggregate $\volR$ diagnostic, or does
cross-cultural variation fragment the diagnostic?

\item \label{oq:m5_interaction}
\textbf{Interaction with axiom M5 (non-triviality).}
Under the capped-benchmark discipline
(Definition~\ref{def:polarity_benchmark}), sufficiency-met needs
contribute a fixed cap.  The weight of the cap is a design
parameter that affects the relative importance of
need-satisfaction versus want-pursuit in $\volL$.

\item \label{oq:need_identification}
\textbf{Need identification.}
How does the framework identify which capabilities are needs
versus wants?  The poset-structural approach (high fan-out = need)
is a good heuristic but not infallible.  A hybrid of structural
position and revealed-sacrifice polarity transition is probably
required.
\end{enumerate}
